% Standalone source-record documentation; no research attempts.
\documentclass[11pt]{article}
\usepackage[margin=25mm]{geometry}
\usepackage{fontspec,unicode-math}
\setmainfont{Latin Modern Roman}
\setmathfont{Latin Modern Math}
\usepackage{amsmath,amssymb,mathtools,xurl,hyperref}
\hypersetup{colorlinks=true,urlcolor=blue,linkcolor=blue}
\setcounter{secnumdepth}{0}
\setlength{\parindent}{0pt}
\setlength{\parskip}{5pt}
\setlength{\emergencystretch}{3em}
\newcommand{\sref}[2]{\hyperlink{source:#1:#2}{[S#2]}}
\newcommand{\cataloguescope}[1]{\paragraph{Source record.}#1}
\newcommand{\floor}[1]{\left\lfloor#1\right\rfloor}
\begin{document}
% UnsolvedMath ID 30006187; source catalogue code OWR-14299084-010
\setcounter{section}{30006186}
\section{Cryptographic Equivalences for High-Energy Phenomena}\label{30006187}
{\small\sffamily UnsolvedMath ID 30006187 · Computer Science}
\subsection{Definitions and mathematical statement}

\paragraph{Shared notation.}Standard number systems are $\mathbb N,\mathbb Z,\mathbb Q,\mathbb R,\mathbb C$. Specialized notation and index ranges follow the cited source and the definitions below.


A cryptographic primitive specifies a task and security experiment against an adversary. A hardness equivalence requires explicit reductions in both directions between computational problems.
\paragraph{Record scope and qualifications.}
The catalogue title is retained as metadata. Named certificates, contracts, bounds or frameworks in that title are not asserted as results of the original source extract.
The source asks for new examples and does not provide a particular phenomenon or equivalence to formulate.

Source-record limitation: Further examples of physics/cryptography equivalences are requested without a mathematical proposition.
\paragraph{Statement recorded in the supplied source.}
Find further high-energy-physics phenomena whose computational hardness is equivalent to the existence of standard cryptographic primitives, extending the known equivalence for decoding Hawking radiation. \sref{30006187}{2}
\paragraph{Source status.}
The supplied catalogue labels this source record open; that metadata was checked against the supplied snapshot on 2026-10-11. The label is preserved for provenance and is not an independent verification of present open status or mathematical correctness. The catalogue literature assessment is dated 2026-08-22.
\subsection{Short English statement}
Find further links between computational physics and cryptographic hardness.
\subsection{Sources}
\begin{itemize}
\item[\hypertarget{source:30006187:1}{S1}] Ulam AI and the UnsolvedMath Contributors, UnsolvedMath: A Curated Collection of Open Mathematics Problems, record 30006187 (OWR-14299084-010). Catalogue curation and metadata: CC BY 4.0. The source extract is retained; definitions, English summaries and qualifications are editorial additions.
\url{https://huggingface.co/datasets/ulamai/UnsolvedMath}.
\item[\hypertarget{source:30006187:2}{S2}] Cryptography (2025), Oberwolfach Reports, DOI 10.4171/owr/2025/3, pp. 141--174
\url{https://doi.org/10.4171/owr/2025/3}.
\item[\hypertarget{source:30006187:3}{S3}] Brakerski, Black-Hole Radiation Decoding is Quantum Cryptography
\url{https://arxiv.org/abs/2211.05491}.
\item[\hypertarget{source:30006187:4}{S4}] OWR 2025/3 Cryptography report
\url{https://ems.press/content/serial-article-files/51349}.
\item Source-specific Oberwolfach Reports licences, including CC BY-SA 4.0 where applicable, are retained; see THIRD\_PARTY\_NOTICES.md.
\end{itemize}
\paragraph{Preparation and review.}
The original source and UnsolvedMath contributors retain attribution for the source extract. Codex (GPT-6; exact serving revision not exposed) assisted with definitions, English summaries, scope notes and formatting. No new human mathematical authorship or independently verified solution is asserted. Permanent source ID 30006187 and code OWR-14299084-010 checked against the complete supplied registry and current repository. Distinct upstream records, including repeated questions and broad or damaged extracts, are retained for complete coverage of the source CS category. A repeated record is not a new independent conjecture. Complete mathematical-equivalence and current-literature review has not been performed. Preparation and final review follow contribution prompts C and D, with source gaps disclosed. Structural and rendering checks do not establish a new theorem or open status.
\label{end:30006187}
\end{document}
